% !PRINTED 152
\weberpage{152}
%<<CONT
formation unchanged and is called an \gesperrt{invariant}. There are two fundamental ones,
which we denote by $g_2$, $g_3$:
\weq{9}{\begin{aligned}
g_2 &= -4(e_2e_3+e_3e_1+e_1e_2) = 2(e_1^2+e_2^2+e_3^2) \\
&= \frac{64}{3}\frac{K^4}{\omega_1^4}(1-\varkappa^2\varkappa'^2),
\end{aligned}}
\weq{10}{g_3 = 4e_1e_2e_3 = \frac{2^8}{27}\frac{K^6}{\omega_1^6}(2+\varkappa^2\varkappa'^2)(\varkappa'^2-\varkappa^2),}
and we add besides the function known under the name of the \gesperrt{discriminant}:
\weq{11}{\begin{aligned}
G &= (e_2-e_3)^2(e_3-e_1)^2(e_1-e_2)^2 = \frac{1}{16}(g_2^3-27g_3^2) \\
&= \frac{(2K)^{12}}{\omega_1^{12}}\varkappa^4\varkappa'^4,
\end{aligned}}
for which, by §\,43 (3), (15) and §\,34 (2), we may also put:
\weq{12}{G = \frac{2^8\pi^{12}\eta(\omega)^{24}}{\omega_1^{12}}.}
$\wp(u)$, $g_2$, $g_3$, $G$ are homogeneous functions, as the following relations show, in which
$\lambda$ is an arbitrary factor.
\weq{13}{\begin{aligned}
\wp(\lambda u, \lambda\omega_1, \lambda\omega_2) &= \lambda^{-2}\wp(u, \omega_1, \omega_2), \\
\wp'(\lambda u, \lambda\omega_1, \lambda\omega_2) &= \lambda^{-3}\wp'(u, \omega_1, \omega_2), \\
g_2(\lambda\omega_1, \lambda\omega_2) &= \lambda^{-4}g_2(\omega_1, \omega_2), \\
g_3(\lambda\omega_1, \lambda\omega_2) &= \lambda^{-6}g_3(\omega_1, \omega_2), \\
G(\lambda\omega_1, \lambda\omega_2) &= \lambda^{-12}G(\omega_1, \omega_2).
\end{aligned}}
If we put, as in §\,45 (2)
\[ \omega_1 = 2K, \qquad \omega_2 = 2iK', \]
we then have
\weq{14}{\begin{aligned}
e_1 &= \frac{1+\varkappa'^2}{3}, \quad e_2 = -\frac{\varkappa'^2-\varkappa^2}{3}, \quad e_3 = -\frac{1+\varkappa^2}{3}, \\
g_2 &= \frac{4}{3}(1-\varkappa^2\varkappa'^2), \quad g_3 = \frac{4}{27}(\varkappa'^2-\varkappa^2)(2+\varkappa^2\varkappa'^2), \\
G &= \varkappa^4\varkappa'^4.
\end{aligned}}
If we eliminate $\omega_1$ from the general expressions (9), (10), (11) for $g_2$, $g_3$, $G$,
we obtain homogeneous functions of order zero, that is, functions of $\omega$ alone, which remain
\gesperrt{unchanged under linear transformations}. Such functions are $g_2^3:G$, $g_3^2:G$. Among these
functions, which can all be reduced to one another, we single out one, which we denote by
%CONT>>
